\chapter{Colors}

\blindtext

\begin{figure}[h]
    \centering
    \begin{tikzpicture}
        \foreach[count=\i] \colour in {uulm-med,uulm-mawi,uulm-nawi,uulm-in,uulm-akzent,uulm} {
            \pgfmathtruncatemacro\j{\i-1}
            \pgfmathparse{\i-0.5}
            \let\textpos\pgfmathresult
                \node[text width=2.5cm, align=right] at (-2,.5*\textpos) {\colour};
            \foreach \x in {0,1,...,9} {
                \pgfmathtruncatemacro\y{\x+1}
                \pgfmathparse{100-\x*10}
                \let\saturation\pgfmathresult
                    \fill[\colour!\saturation] (.5*\x,.5*\j) rectangle (.5*\y,.5*\i);
            }
        }
    \end{tikzpicture}
    \caption{The Ulm University Colour Palette\cite{cd}}
\end{figure}

\blindtext

\begin{table}
    \centering
    \renewcommand{\arraystretch}{1.3}
    \scriptsize
    \begin{tabular}{lp{3.2cm}lll}
        \textbf{Color Macro} & \textbf{Association}  & \textbf{CMYK} & \textbf{sRGB} & \textbf{Hex} \\\midrule 
        uulm & Ulm University  & C30-mo-yo-k35 & 125-154-170 & \#7D9AAA \\
        uulm-akzent & Emphasis Colour  & C0-mo-y28-k38 & 169-162-141 & \#A9A280\\
        uulm-in & Engineering, Computer\newline Science, and Psychology & C0-m100-y63-k29 & 163-38-56 & \#A32638 \\
        uulm-nawi & Natural Sciences & C0-m66-y100-k7 & 189-96-5 & \#BD6005 \\
        uulm-mawi & Mathematics and\newline Economics & C59-mo-y100-k7 & 86-170-28 & \#56AA1C \\
        uulm-med & Medicine & C100-m56-yo-k23 & 38-84-124 & \#26247C\\
    \end{tabular}
    \caption{Color Definition of the Ulm University Colors~\cite{cd}}
\end{table}
